\section{Conclusion}
\label{sec:conclusion}

We presented a linguistically informed, morphology-first hybrid tokenizer designed for Turkish and similar agglutinative languages. The tokenizer combines curated root and affix lexicons with phonological normalization (mapping surface allomorphs to shared identifiers) and a controlled subword fallback for coverage. This design aims to produce token sequences that more closely align with morpheme boundaries while remaining practical for large-scale NLP pipelines.

On TR-MMLU, the proposed tokenizer achieves 90.29\% TR~\% and 85.80\% Pure~\%, indicating substantially stronger morpheme-level alignment than several general-purpose tokenizers. We additionally report downstream sentence embedding evaluation on STS and TR-MTEB using \textbf{randomly initialized} models to isolate tokenizer effects from pretrained knowledge. The TurkishTokenizer-based model reaches \textbf{51.44\%} Pearson correlation on STSb-TR, compared to 43.01\% for the Tabi baseline-a gain of \textbf{+8.43 percentage points}. On TR-MTEB, TurkishTokenizer achieves 39.57\% overall average compared to 34.92\% for Tabi (+4.65 points). These gaps demonstrate that morphology-first tokenization provides a stronger inductive bias for learning Turkish semantic representations from scratch, and the same tokenizer also yields the strongest average accuracy on a centroid-based TurBLiMP minimal-pairs proxy.

We emphasize that empirical claims in this paper are Turkish-focused. We outline concrete next steps-improved morphophonological handling, better capitalization edge cases, and standardized efficiency measurements-in Section~\ref{sec:future_work}.
